\begin{definition}
	\textit{Ступенчатой} (кусочно-постоянной) функцией на измеримом множестве $E\subset\R^n$ называется измеримая функция, множество значений которой конечно или счётно.
\end{definition}
\begin{theorem}\label{24_th4}
	Если $f:E\ra\R$ измерима, то найдётся возрастающая (убывающая) последовательность ступенчатых на $E$ функций, сходящаяся к $f$ на $E$.
\end{theorem}
\begin{proof}
	Разобьём $\R$ на полуинтервалы:
	\[\R=\bsc_{l\in\Z}\l[\frac{l}{2^m},\frac{l+1}{2^m}\r),\ 
	E_{l,m}:=\l\{x\in E:\frac l{2^m}\le x<\frac{l+1}{2^m}\r\}=E_f\l(\frac{l+1}{2^m}\r)\sm E_f\l(\frac l{2^m}\r)\]
	\[\bsc_{l\in\Z}E_{l,m}=E\]
	Пусть $g_m(x):=\frac l{2^m},\ x\in E_{l,m}$, тогда:
	\[(\fa m\in\N\cup\{0\})(\fa x\in E)\ g_m(x)\le g_{m+1}(x)\]
	\[(\fa x\in E)(\fa m\in\N\cup\{0\})\ 0\le f(x)-g_m(x)<\frac1{2^m}\]
	Следовательно, $g_m\os{E}{\rra}f$. Получили требуемое.
\end{proof}
\begin{definition}
	Последовательность измеримых на измеримом множестве $E\subset\R^n$ функций $f_m$ \textit{сходится по мере} к $f$ на $E$, если:
	\[(\fa\eps>0)\ \lim_{m\ra\i}\mu\{x\in E:|f(x)-f_m(x)|\ge\eps\}=0\]
	Обозначается как $f_m\Ra f$
\end{definition}
\begin{definition}
	\textit{Почти всюду} - всюду, кроме, может быть, множества, лебегова мера которого равна 0.
\end{definition}
\begin{theorem}\nf{(Связь сходимости по мере и почти всюду)}\label{24_th5}\\
	Если последовательность измеримых и почти всюду конечных на измеримом множестве $E$ конечной меры функций $f_m$ сходится почти всюду на $E$ к почти всюду конечной функции $f$, то $f_m$ сходится к $f$ по мере.
\end{theorem}
\begin{proof}
	Пусть $E_0$ - множество точек множества $E$, в которых хотя бы одна из функций $f_m$ равна $\pm\i$, или $f$ равна $\pm\i$, или $f_m\nra f(x)$.\\ Заметим, что:
	\[\mu(E_0)=0\]
	Положим: $\{\eps_k\}$, $\lim_{k\ra\i}\eps_k=0$, $\eps_k$ убывает, тогда:
	\[E_m(\eps_k):=\{x\in E:|f(x)-f_m(x)|\ge\eps\}\]
	Докажем, что $\ds\bc_{k=1}^\i\vlim[m\ra\i]E_m(\eps_k)\subset E_0$\\
	Пусть $\ds x\notin\bc_{k=1}^\i\vlim E_m(\eps_k)$, что может произойти тогда и только тогда, когда:
	\[(\fa k\in\N)(\ex N\in\N)(\fa m>N)\ x\notin E_m\lra|f(x)-f_m(x)|<\eps\]
	То есть имеем: $f_m(x)\us{m\ra\i}{\ra}f(x)\lra x\notin E_0$.\\
	Теперь воспользуемся (\ref{21_l3}):
	\[\vlim\mu(E_m)\le\mu\l(\vlim E_m\r)\le\mu(E_0)=0\]
	Получили требуемое.
\end{proof}
\begin{note}
	Обратное утверждение к (\ref{24_th5}) неверно. Покажем на следующем примере.
\end{note}
\begin{example}\nf{Пример Ференца Рисса}\\
	\[(\fa m\in\N)\ m=2^k+l,\ l=0,\dots,2^k-1,\ f_m(x):=1_{\l[\frac l{2^k},\frac{l+1}{2^k}\r]}(x),\ x\in[0,1]\]
	Рассмотрим множества $E_m$:
	\[E_m=\{x\in[0,1]:|f_m(x)|\ge\eps\},\ \mu(E_m)=\frac1{2^k}\us{m\ra\i}{\ra}0\]
	Тогда, так как всегда можно найти подпоследовательность функций, значение которых в точке $x$ равны 1:
	\[(\fa x\in[0,1])\ f_m(x)\us{m\ra\i}\nra f(x)\]
\end{example}
\begin{note}
	Для множеств бесконечной меры утверждение неверно.\\
	Пусть $E=\R$,\ $f_n(x)=1_{[-m,m]}$, тогда:
	\[(\fa x\in\R)\ \lim_{m\ra\i}f_m(x)=1=f(x)\]
	\[(\fa\eps>0\text{ и }\eps\le1)\ \mu\{x\in\R:|f_m(x)-f(x)|\ge\eps\}=+\i\]
\end{note}
\begin{example}
	$1_A(x)$ измерим тогда и только тогда, когда $A$ измеримо по Лебегу.
	\[E_{1_A}(a)=
	\begin{cases}
		\R^n,\ a>1\\
		\R^n\sm A,\ 0<a\le1\\
		\emptyset,\ a\le0
	\end{cases}
	\]
	Для того, чтобы $E_{1_A}$ было измеримым, нужно чтобы $\R^n\sm A$ было измеримо, что равносильно измеримости A.\\
	Получили требуемое.
\end{example}
\begin{note}
	Для любого измеримого множества $B$, $\mu(B)>0$ существует неизмеримое по Лебегу подмножество $E\subset B$.
\end{note}
\begin{proof}
	Для $B\subset[0,1]$ доказательство практически не отличается от доказательства существования неизмеримого множества.
\end{proof}
\begin{example}
	Канторова лестница. Вспомним канторово множество:
	\[P=\bigcap_{m=1}^\i I_m,\ I_{m+1}=I_m\sm J_{m+1}\]
	Концы выбрасываемых интервалов из $J_m$ образуют множество $\ds E=\l\{\frac{k}{3^m},\ k=0,\dots,3^m,\ \m=0,1,\dots\r\}$\\
	Определим $\v\phi(x)$ на множестве $E$:
	\[\v\phi(0):=0,\ \v\phi(1):=1\]
	На концах интервалов определим $\v\phi$ как среднее арифметическое из ближайших с разных сторон уже определённых значений.\\
	Множество значений $\v\phi$: $F=\l\{\frac k{2^m},\ k=0,\dots,2^m,\ m=0,1,\dots\r\}$\\
	Теперь определим канторову лестницу:
	\[\phi(x):=\sup_{y\in E:y\le x}\v\phi(y)\]
	Заметим, $\phi(x)$ возрастает.\\
	Докажем, что $\phi(x)$ непрерывна. От противного, пусть $x_0\in[0,1]$ - точка разрыва первого рода, тогда на интервале $(\phi(x_0-0),\phi(x_0+0))$ нет значений $\phi$, что невозможно, так как $F\subset\phi([0,1])$
\end{example}
\begin{example}
	Построим неизмеримую функцию, являющуюся композицией измеримой и непрерывной функций.\\
	Пусть $\psi(X):=\frac12(\phi(x)+x)$, где $\phi$ - канторова лестница.\\
	Заметим, $\psi$ - непрерывна и строго возрастает на $[0,1]$, $\psi([0,1])=[0,1]$
	\[\mu(\psi([0,1]\sm P))=\frac12\Ra\mu(\psi(P))=1-\mu(\psi([0,1]\sm P))=\frac12\]
	\[\psi(P)\text{ - имеет меру}>0\Ra\ex B\subset\psi(P),\ B\text{ - неизмеримое по Лебегу}\]
	По теореме об обратной функции $\ex f=\psi^{-1}$ на $[0,1]$, при этом, $\psi$ строго возрастает.
	Положим:
	\[A:=f(B)\subset P,\ \mu(P)=0\Ra\mu(A)=0\]
	Следовательно, $1_A$ - измеримая функция, при этом $1_A\circ f=1_B$, так как $B$ - неизмеримо, то требуемое получено.
\end{example}